\section{RLHF：把人类比较转化为策略更新}

本章从 slide 上的两步图出发，恢复完整的 SFT--RM--RL 教学链。核心思想是：人类不必为每个输出写一个绝对分数，而只需比较少量候选；reward model 学习这种排序规律，再为大量未人工标注的样本提供可优化的近似信号。

\subsection{第一步：从 pairwise comparison 学 reward model}

给定 prompt $x$，策略产生两个回答 $y_w$ 与 $y_l$，标注者选择更符合意图的 winner。reward model $r_\phi(x,y)$ 不直接生成文本，而是把 prompt--response 对映射为标量。常用 Bradley--Terry 形式把奖励差转成偏好概率：

\[
P_\phi(y_w \succ y_l\mid x)
=\sigma\!\left(r_\phi(x,y_w)-r_\phi(x,y_l)\right),
\]

其中，$\sigma(z)=1/(1+e^{-z})$ 是 logistic 函数；$r_\phi$ 是参数为 $\phi$ 的 reward model；符号 $y_w \succ y_l$ 表示标注者更偏好 $y_w$。对应的负对数似然为

\[
\mathcal{L}_{\mathrm{RM}}(\phi)
=-\mathbb{E}_{(x,y_w,y_l)\sim D}
\log \sigma\!\left(r_\phi(x,y_w)-r_\phi(x,y_l)\right).
\]

其中，$D$ 表示人类比较数据集，期望对 prompt 与成对回答取平均。损失只约束相对次序，因此奖励的绝对零点没有独立意义。

\lecturefigure{slide-06-rlhf-reward-model.jpg}{RLHF 的第一步：由人类比较训练 reward model。}{Stanford Online 官方视频 00:06:35--00:07:55。}

\begin{knowledgebox}{读图：比较数据解决什么问题}
图中人类面对同一 prompt 的多个 response，提供相对排序；reward model 学会预测这种排序。相对判断通常比写出精确奖励更自然，也把一次人工比较摊薄到许多后续 policy samples 上。它没有消除主观性：不同 labeler 会不一致，最终模型学到的是数据分布中的聚合偏好，而不是客观真理的直接读数。
\end{knowledgebox}

\teachervoice{讲者明确说，labeler 经常并不一致，模型会在训练中聚合这些差异。这里的工程问题不只是 ``标签噪声''，还包括谁被雇来标注、给他们什么规范、哪些群体没有进入数据，以及冲突偏好是否应该被平均。}

\subsection{第二步：优化 policy，同时限制偏移}

获得 reward model 后，可以采样策略回答并用强化学习提高预测奖励。若只最大化 $r_\phi$，policy 可能钻 reward model 的漏洞；InstructGPT 因而使用对参考策略的 KL 约束，并混入 pretraining loss 以降低能力回退。一个教学化目标是

\[
J(\theta)=
\mathbb{E}_{x\sim D_x,\,y\sim\pi_\theta(\cdot\mid x)}
\left[r_\phi(x,y)-\beta\,
D_{\mathrm{KL}}\!\left(\pi_\theta(\cdot\mid x)\,\|\,\pi_{\mathrm{ref}}(\cdot\mid x)\right)\right]
+\gamma\,\mathbb{E}_{z\sim D_{\mathrm{ptx}}}\log \pi_\theta(z).
\]

其中，$\pi_\theta$ 是待优化 policy，$\pi_{\mathrm{ref}}$ 是 SFT/reference policy，$\beta$ 控制偏离参考模型的惩罚强度，$D_{\mathrm{KL}}$ 衡量两个输出分布的差异，$D_x$ 是 prompt 分布，$D_{\mathrm{ptx}}$ 是预训练文本，$\gamma$ 控制保留语言建模能力的权重。

\lecturefigure{slide-07-rlhf-full-loop.jpg}{完整的 slide 叙事：reward model 提供反馈，预训练语言模型经 RL 得到新 policy。}{Stanford Online 官方视频 00:07:55--00:09:24。}

\begin{importantbox}{课堂两步图与论文三阶段并不矛盾}
Slide 为了突出 reward modeling，把流程压成 ``比较训练 RM'' 与 ``RL 优化 policy''。讲者口头提醒还有一步被省略。按 InstructGPT 论文，常见流程是：先用 demonstrations 做 supervised fine-tuning（SFT），再用 comparisons 训练 reward model（RM），最后用 PPO 等 RL 方法优化 policy。讲义补回 SFT 是来源澄清，不是假装 slide 上出现了第三个框。
\end{importantbox}

\begin{lstlisting}[caption={RLHF 三阶段的最小伪代码}]
policy = supervised_finetune(base_model, demonstrations)
reward_model = train_pairwise_reward_model(policy, comparisons)

for batch in prompt_batches:
    responses = policy.sample(batch)
    rewards = reward_model.score(batch, responses)
    policy = ppo_update(
        policy,
        rewards=rewards,
        reference=supervised_policy,
        kl_penalty=beta,
        pretraining_batch=next(pretraining_data),
    )
\end{lstlisting}

\subsection{为什么不是只做 imitation learning}

SFT/behavioral cloning 学习 ``人类会怎样回答''；preference optimization 学习 ``人类更喜欢哪个结果''。讲者指出一个双向限制：某些任务人类比模型强，模型可能无法完美模仿；另一些任务模型已有不同优势，强制模仿人类动作路径反而限制它。RL 允许模型寻找自己的实现方式，只要结果通过偏好评价。

\begin{knowledgebox}{术语消化：SFT、RM、PPO、PPO-ptx}
\begin{tabularx}{\textwidth}{L{2.2cm} >{\raggedright\arraybackslash}X >{\raggedright\arraybackslash}X}
\toprule
术语 & 解决的问题 & 在本讲中的位置 \\
\midrule
SFT & 用高质量 demonstrations 建立可交互的初始 policy & 提供参考策略与基本指令跟随能力 \\
RM & 把成对偏好压缩为可重复调用的标量函数 & 让有限人类比较监督大量样本 \\
PPO & 在限制单次更新幅度的情况下提高预测奖励 & 是当时熟悉且有效的优化器，不是已证明最优 \\
PPO-ptx & RL 时混入预训练目标 & 缓解只追偏好奖励带来的能力回退 \\
\bottomrule
\end{tabularx}
\end{knowledgebox}

\teachervoice{学生问为什么用 PPO。讲者的回答非常工程化：团队熟悉它，而且 ``works pretty well''。这说明算法选择包含历史路径与工具成熟度，不能从 ``被采用'' 推出 ``理论上最佳''。}

\subsection{本章小结}

RLHF 用人类比较训练 reward model，再优化 policy；完整实践还包含 SFT、KL regularization 与预训练数据保留。它把反馈成本从每次动作的人类评分转成可复用模型，但也把 labeler 偏好、reward hacking 和 proxy misspecification 带入训练目标。

